\documentclass[10pt]{article}
\usepackage{amsmath,amssymb,amsthm}
\usepackage[T1]{fontenc}
\usepackage{lmodern}
\usepackage{hyperref}
\allowdisplaybreaks
\newtheorem{theorem}{Theorem}[section]
\newtheorem{proposition}[theorem]{Proposition}
\newtheorem{lemma}[theorem]{Lemma}
\theoremstyle{remark}
\newtheorem{remark}[theorem]{Remark}
\newcommand{\N}{\mathbb N}
\newcommand{\E}{\mathbb E}
\newcommand{\Prob}{\mathbb P}
\newcommand{\ind}{\mathbf 1}
\newcommand{\dd}{\,dt}
\newcommand{\dens}{\mathop{\mathrm{dens}}}
\newcommand{\cn}{c_{\mathrm N}}
\newcommand{\cd}{c_{\mathrm D}}
\title{The density interpretation of Theorem 4 of Erd\H{o}s, Graham, Ruzsa and Straus}
\author{Haoyu Chen\\ Independent researcher}
\date{}
\begin{document}
\maketitle

\begin{abstract}
For each fixed $0<\alpha\leq 1/2$, we give the natural-density limits of
the proportions of integers up to $n^\alpha$ which divide, or do not
divide, $\binom{2n}{n}$. The constants have absolutely convergent integral
series and effective truncation bounds. The nondivisor constant tends to
zero as $\alpha\downarrow0$, whereas the divisor constant tends to one.
Neither proportion has a limit along all integers $n$ in this range.
We prove the uniform prime counts needed for the localized moment
calculation and give certified numerical intervals. No claim
is made for $1/2<\alpha<1$, or for an effective exceptional-set rate in $n$.
\end{abstract}

\section{Introduction}
Write $B_n=\binom{2n}{n}$. Theorem 4 on p.~89 of
Erd\H{o}s--Graham--Ruzsa--Straus \cite{EGRS} reads as follows
(the mathematical notation and the final parenthesis are retained):
\begin{quote}
\textsc{Theorem 4.} \textit{For $\alpha<1$,}
\[
 \left|\left\{m:1\leq m\leq n^\alpha,\ m\nmid\binom{2n}{n}\right\}\right|
 =c(\alpha)n^\alpha+o(n^\alpha),
\]
\textit{where $c(\alpha)\to1$ as $\alpha\to0$.
(In fact, $c(\alpha)$ can be explicitly calculated.)}
\end{quote}
The qualifier ``for almost all $n$'' occurs in the preceding display (6)
and its introduction, but not in the theorem itself. In the range
$0<\alpha\leq1/2$, an interpretation as a limit along all $n$ is false.
Even after the density qualifier is supplied, the direction of
divisibility in the assertion about the constant tending to one is
reversed: it is the divisor proportion that has that limit.

We prove the density statement for $0<\alpha\leq1/2$, with an explicit
series for both constants and certified intervals at four values of
$\alpha$. We also correct the endpoint interpretation of display (6).
Nothing here is asserted for $1/2<\alpha<1$, although the printed theorem
states the wider range $\alpha<1$.

\section{Related work}
The main question of \cite{EGRS}, whether
$f(n)=\sum_{p\leq n,\ p\nmid B_n}1/p$ is bounded, is Problem~377 in
Bloom's list of Erd\H{o}s problems \cite{EP377}; this note does not
address it. Several later papers determine densities of sets of
\emph{indices} $n$ defined by divisibility properties of $B_n$.
Pomerance \cite{Po15} shows that, for each fixed $k\geq1$, $n+k\mid B_n$
for a set of $n$ of density one. Sanna \cite{Sa18} bounds the upper
density of $\{n:n\mid B_n\}$. Ford and Konyagin \cite{FK21} prove that
$\{n:n^\ell\mid B_n\}$ has a positive density for each $\ell\geq1$ and
evaluate these densities numerically. Pomerance \cite{Po26} proves
density-one divisibility statements for products of consecutive
integers. These results concern the density of indices $n$, not the
proportion of divisors of $B_n$ among $m\leq n^\alpha$ for a typical
$n$, which is the quantity of Theorem~4 of \cite{EGRS} and of this note.
In the formal-conjectures repository, a change merged in September 2026
\cite{FC5792} replaced a cofinite-convergence formalization of the
reciprocal-prime result of \cite{EGRS} by a density-one statement; that
is a different statement from Theorem~4.

In a bounded search completed on 10 October 2026 (reviews and forward
citations of \cite{EGRS}, the problem page and its forum, OEIS, and
public code repositories), we did not locate an earlier published
correction of the printed Theorem~4, or an earlier explicit evaluation
of the densities studied here. The MathSciNet review of \cite{EGRS} was
not accessible to us.

\section{Statement and constants}\label{sec:statement}
For real $x=n^\alpha$ put
\[
 N_\alpha(n)=\frac{\#\{1\leq m\leq x:m\nmid B_n\}}x,
 \qquad
 D_\alpha(n)=\frac{\#\{1\leq m\leq x:m\mid B_n\}}x.
\]
Thus $N_\alpha(n)+D_\alpha(n)=\lfloor x\rfloor/x$, rather than exactly
one. Convergence in natural density to $c$ means that, for every $\eta>0$,
$X^{-1}\#\{n\leq X:|u(n)-c|>\eta\}\to0$. Throughout the proof $X$ may
be restricted to positive integers; replacing a real $X$ by $\lfloor
X\rfloor$ changes normalized counts by at most $1/X$.

For $0<t<1$ define
\[
 w(t)=2^{-\lfloor1/t\rfloor},\qquad g(t)=\frac{w(t)}t,
 \qquad A(y)=\int_0^y g(t)\,dt\quad(0<y\leq1).
\]
Values at the countably many breakpoints do not affect integrals. Set
\[
 J_0(\alpha)=1,\qquad
 J_k(\alpha)=\frac1{k!}
 \int_{\substack{t_i>0\ (1\leq i\leq k)\\t_1+\cdots+t_k\leq\alpha}}
       \prod_{i=1}^k g(t_i)\,dt_1\cdots dt_k\quad(k\geq1).
\]

\begin{theorem}[Corrected Theorem 4, for $0<\alpha\leq1/2$]
\label{thm:main}
For every fixed $0<\alpha\leq1/2$ and every $\eta>0$,
\begin{align*}
 \lim_{X\to\infty}\frac1X\#\{n\leq X:|N_\alpha(n)-\cn(\alpha)|>\eta\}&=0,\\
 \lim_{X\to\infty}\frac1X\#\{n\leq X:|D_\alpha(n)-\cd(\alpha)|>\eta\}&=0,
\end{align*}
where the absolutely convergent series are
\begin{equation}\label{eq:density}\tag{DENSITY}
 \begin{aligned}
 \cd(\alpha)&=\sum_{k=0}^{\infty}(-1)^kJ_k(\alpha),\\
 \cn(\alpha)&=\sum_{k=1}^{\infty}(-1)^{k+1}J_k(\alpha)=1-\cd(\alpha).
 \end{aligned}
\end{equation}
If $s=\lfloor1/\alpha\rfloor\geq2$, then
\begin{align}
 A(\alpha)&=2^{-s}\log((s+1)\alpha)
       +\sum_{r=s+1}^{\infty}2^{-r}\log(1+1/r),\label{eq:A}\\
 0<A(\alpha)&\leq A_s:=\frac{2^{1-s}}s\leq\frac14,
 \qquad 0\leq J_k(\alpha)\leq\frac{A(\alpha)^k}{k!},\label{eq:bounds}\\
 A(\alpha)(1-A(\alpha)/2)&\leq\cn(\alpha)\leq A(\alpha)<\frac14.
 \label{eq:bracket}
\end{align}
For every integer $K\geq0$, writing
$E_K(u)=u^{K+1}/((K+1)!(1-u/(K+2)))$,
\begin{equation}\label{eq:tail}\tag{TAIL}
 \left|\cd(\alpha)-\sum_{k=0}^K(-1)^kJ_k(\alpha)\right|
 \leq E_K(A(\alpha))\leq E_K(A_s).
\end{equation}
The same bound holds for the nondivisor series truncated at $k=K$
(the empty sum is zero). Both proportions also converge in Ces\`aro mean.
\end{theorem}

\begin{proposition}\label{prop:nolimit}
For every fixed $0<\alpha\leq1/2$, neither $N_\alpha(n)$ nor
$D_\alpha(n)$ has a limit as $n\to\infty$ through all integers.
For every integer $h\geq0$,
\[
 v_2\left(\binom{2^{h+1}}{2^h}\right)=1,\qquad
 N_\alpha(2^h)\geq\frac{\lfloor2^{h\alpha}/4\rfloor}{2^{h\alpha}}
 \geq\frac14-2^{-h\alpha}.
\]
Moreover $\cn(\alpha)\to0$ and $\cd(\alpha)\to1$ as $\alpha\downarrow0$.
\end{proposition}

\begin{remark}[Effectiveness and scope]\label{rem:scope}
The constant can be enclosed to any prescribed rational accuracy when
$\alpha$ is rational, and for a computable real $\alpha$ when an
approximation procedure for $\alpha$ is supplied. This is a claim about
the construction in Section~\ref{sec:numerics}; the supplied program
implements $\alpha=1/s$ for integers $s\geq2$. All numerical error
bounds below are explicit. We give no effective onset or rate for the
exceptional sets in $n$. The quantifiers in Theorem~\ref{thm:main} are for
each fixed $\alpha$ and $\eta$, not uniformly in them. Nothing is claimed
for $1/2<\alpha<1$.
\end{remark}

\begin{lemma}[Integral bounds and alternating-series bracket]
\label{lem:alternating}
The formulas and inequalities \eqref{eq:A}--\eqref{eq:tail} hold for
the constants defined by \eqref{eq:density}.
\end{lemma}
\begin{proof}
Integrate $2^{-r}/t$ on each interval $(1/(r+1),1/r)$ to obtain
\eqref{eq:A}. Since $\log(1+1/r)\leq1/r$ and $\alpha\leq1/s$,
\[
 A(\alpha)\leq\sum_{r\geq s}\frac{2^{-r}}r
 \leq\frac1s\sum_{r\geq s}2^{-r}=A_s.
\]
Dropping the simplex constraint gives $J_k\leq A(\alpha)^k/k!$.
More particularly, integrate out the last coordinate and then enlarge
its permitted interval to $(0,\alpha)$ to obtain
\begin{equation}\label{eq:successive}
 (k+1)!J_{k+1}(\alpha)
 \leq k!J_k(\alpha)A(\alpha),\qquad
 J_{k+1}(\alpha)\leq\frac{A(\alpha)}{k+1}J_k(\alpha).
\end{equation}
Thus the nonnegative terms decrease to zero, and the alternating-series
bracket gives $J_1-J_2\leq\cn\leq J_1$. Since $J_1=A$ and $J_2\leq
A^2/2$, this proves the non-strict parts of \eqref{eq:bracket}.
The strict inequality follows from
\[
 A(\alpha)\leq A(1/2)=\sum_{r\geq2}2^{-r}\log(1+1/r)
 <\frac12\sum_{r\geq2}2^{-r}=\frac14.
\]
Finally, for $j\geq0$ and $0\leq u\leq1/4$,
\[
 \frac{u^{K+1+j}}{(K+1+j)!}
 \leq\frac{u^{K+1}}{(K+1)!}\left(\frac{u}{K+2}\right)^j.
\]
The geometric sum bounds the absolute tail. Its terms increase with
$u$, so replacing $A(\alpha)$ by $A_s$ is permitted.
\end{proof}

\section{Proof of the density theorem}\label{sec:proof}
\subsection{Sieve reduction}
For $0<\delta<\alpha$, let $\mathcal P_\delta(n)$ be the primes
$n^\delta<p\leq n^\alpha$ with $p\nmid B_n$. Let $F_\delta(n)$ be
$n^{-\alpha}$ times the number of integers $1\leq m\leq n^\alpha$
having no prime factor in $\mathcal P_\delta(n)$, and put
\[
 U_\delta(n)=\frac{\lfloor n^\alpha\rfloor}{n^\alpha}-F_\delta(n),
 \quad
 G_\delta(n)=\sum_{\substack{p\leq n^\delta\\p\nmid B_n}}\frac1p,
 \quad
 T(n)=\sum_{\substack{p\leq n^\alpha\\v_p(B_n)\geq1}}
                  p^{-v_p(B_n)-1}.
\]
Every integer counted in $U_\delta$ is a nondivisor. Conversely, if
$m\nmid B_n$, some prime satisfies $v_p(m)>v_p(B_n)$. Either $p\nmid
B_n$, or $p^{v_p(B_n)+1}\mid m$ with $v_p(B_n)\geq1$.
The fraction of multiples of any positive integer $d$ is at most $1/d$.
The union bound therefore proves the pointwise sandwich
\begin{equation}\label{eq:sandwich}
 U_\delta(n)\leq N_\alpha(n)
 \leq U_\delta(n)+G_\delta(n)+T(n).
\end{equation}

\subsection{The localized moment calculation}\label{subsec:local}
We distinguish the printed counts from the uniform bounds proved below.
Write
\[
 \begin{aligned}
 A_1(p;X)&=\#\{p\leq n\leq X:p\nmid B_n\},\\
 A_2(p,q;X)&=\#\{\max(p,q)\leq n\leq X:\gcd(pq,B_n)=1\}.
 \end{aligned}
\]
Changing a strict upper endpoint in a printed count to $n\leq X$
changes its normalized value by at most $1/X$.
On p.~86 of \cite{EGRS}, immediately following (4), the text says
``for each $p$ with $x^{1/(r+1)+\epsilon}<p<x^{1/r-\epsilon}$'',
and the last, unnumbered display is
\begin{equation}\label{eq:input1}
 A(p;x)=x/2^r+o(x).
\end{equation}
Here the printed $A(p;x)$ counts $p\leq n<x$ with $p\nmid B_n$.
For the class-III calculation on p.~88, take
$x^{1/(r+1)}<p\leq x^{1/r}$ and
$x^{1/(t+1)}<q\leq x^{1/t}$, with $p\ne q$, and
order the powers $p^i$ ($1\leq i\leq r$) and $q^j$ ($1\leq j\leq t$)
as $w_1<\cdots<w_{r+t}$. The printed condition (5) is
\[
 \frac{w_{i+1}}{w_i}>x^\epsilon\ (1\leq i<r+t),\qquad
 \frac{x}{w_{r+t}}>x^\epsilon.
\]
The text then says ``then we see by (1) that'', followed by
\begin{equation}\label{eq:input2}
 A(p,q;x)=\frac{x}{2^{r+t}}+o(x).
\end{equation}
Neither quoted display prints a uniformity assertion. The $o(x)$
terms here are part of the quotations; we do not use them as uniform
remainders. We now prove the uniform bounds needed for the prime
summations, independently of statement (1) of \cite{EGRS}.

\begin{lemma}[Counting on nested scales]\label{lem:nested}
Let $X,m$ be positive integers and let $Q\geq2$ be real. Suppose the
positive integers $w_1<\cdots<w_m$ satisfy
\[
 w_1\geq Q,\qquad w_{i+1}/w_i\geq Q\ (1\leq i<m),\qquad X/w_m\geq Q.
\]
If $C(X;w_1,\ldots,w_m)$ counts the integers $1\leq n\leq X$ for
which $n\bmod w_i<w_i/2$ for every $i$, then
\begin{equation}\label{eq:nested}
 \left|C(X;w_1,\ldots,w_m)-2^{-m}X\right|\leq\frac{8mX}{Q}.
\end{equation}
In particular $Q=X^\epsilon$ gives an error at most
$8mX^{1-\epsilon}$, uniformly over every collection of moduli satisfying
the displayed hypotheses.
\end{lemma}
\begin{proof}
Put $h_i(n)=\ind_{\{n\bmod w_i<w_i/2\}}$ and
$c_i=\lceil w_i/2\rceil/w_i$. Each consecutive block of $w_i$
integers contains exactly $c_iw_i$ accepted residues. Splitting any
interval $I$ of consecutive integers into complete blocks and one
remainder of length less than $w_i$ therefore proves
\begin{equation}\label{eq:perioderror}
 \left|\sum_{n\in I}(h_i(n)-c_i)\right|\leq w_i.
\end{equation}
For $i<m$, let $H_{i+1}(n)=\prod_{j=i+1}^m h_j(n)$ on
$\{1,\ldots,X\}$. In this interval $h_j$ changes value at most
$2X/w_j+2$ times, since its only changes in a period occur at residues
$0$ and $\lceil w_j/2\rceil$. Every change of $H_{i+1}$ occurs at
a change of one of its factors. Thus its set of ones is a disjoint
union of at most
\[
 R_i\leq1+\sum_{j=i+1}^m(2X/w_j+2)
 \leq1+3X\sum_{j=i+1}^m\frac1{w_j}
 \leq1+\frac{6X}{w_{i+1}}\leq\frac{7X}{w_{i+1}}
\]
intervals of consecutive integers. The second inequality uses
$X/w_j\geq Q\geq2$, the third uses
$\sum_{j=i+1}^m1/w_j\leq(1/w_{i+1})\sum_{k\geq0}Q^{-k}
\leq2/w_{i+1}$, and the last uses $X/w_{i+1}\geq1$.
Apply \eqref{eq:perioderror} on each such interval to get
\[
 \left|\sum_{n=1}^X(h_i(n)-c_i)H_{i+1}(n)\right|
 \leq R_iw_i\leq7X/Q\qquad(i<m).
\]
For $i=m$ the same sum with the empty product equal to one has
absolute value at most $w_m\leq X/Q$. The exact identity
\[
 \prod_{i=1}^m h_i-\prod_{i=1}^m c_i
 =\sum_{i=1}^m\left(\prod_{j<i}c_j\right)(h_i-c_i)
                         \left(\prod_{j>i}h_j\right)
\]
and $0\leq c_i\leq1$ now give
$|C-X\prod_i c_i|\leq(7(m-1)+1)X/Q\leq7mX/Q$.
Finally, $0\leq c_i-1/2\leq1/(2w_i)$, so another telescoping
identity gives
\[
 0\leq\prod_i c_i-2^{-m}\leq\sum_i(c_i-1/2)
 \leq\frac{m}{2Q}.
\]
The triangle inequality proves \eqref{eq:nested}, since $7+1/2<8$.
\end{proof}

\emph{Application to the prime counts.}
Fix $0<a<b<1$ and $0<\epsilon<a/4$, and suppose that the integer
$X\geq2^{1/\epsilon}$. Put $Q=X^\epsilon$ and $L=\lceil2/a\rceil$.
By the factorial valuation identity \eqref{eq:valuation},
\[
 p\nmid B_n\quad\Longleftrightarrow\quad
 n\bmod p^i<p^i/2\quad\hbox{for every }i\geq1.
\]
Indeed its $i$th summand is one precisely when
$n\bmod p^i\geq p^i/2$, and all its summands are nonnegative.
For an integer $1\leq r\leq L$, if $X^a<p\leq X^b$ and
$X^{1/(r+1)+\epsilon}<p<X^{1/r-\epsilon}$, then
$p^{r+1}>XQ^{r+1}\geq4X$. For $n\leq X$, all conditions with
$i\geq r+1$ are therefore automatic. The remaining moduli are
$p,p^2,\ldots,p^r$; their successive ratios and their smallest modulus
exceed $Q$, and $X/p^r>X^{r\epsilon}\geq Q$.
Lemma~\ref{lem:nested}, followed by deletion of the fewer than $p$
integers with $n<p$, proves the explicit uniform estimate
\begin{equation}\label{eq:uniform1}
 \left|A_1(p;X)/X-2^{-r}\right|
 \leq8rX^{-\epsilon}+p/X\leq8LX^{-\epsilon}+X^{b-1}.
\end{equation}
For two distinct primes $p,q$ in these trimmed bands with indices
$r,t\leq L$, assume also the scale inequalities of (5), with $x=X$.
All powers beyond $p^r$ and $q^t$ again give automatic conditions.
The merged list has $m=r+t\leq2L$ distinct moduli: equality between
a positive power of $p$ and one of $q$ would contradict unique
factorization. Its smallest modulus exceeds $Q$, and (5) supplies
the other hypotheses of Lemma~\ref{lem:nested}. Deleting the fewer
than $\max(p,q)$ initial integers gives
\begin{equation}\label{eq:uniform2}
 \left|A_2(p,q;X)/X-2^{-(r+t)}\right|
 \leq8(r+t)X^{-\epsilon}+\frac{\max(p,q)}X
 \leq16LX^{-\epsilon}+X^{b-1}.
\end{equation}
These estimates prove, with constants independent of the retained
primes, the uniform forms required below.

For the mean of the full sum we use Theorem 2, pp.~86--87:
\begin{equation}\label{eq:inputmean}
 \begin{gathered}
 f(n):=\sum_{\substack{p\leq n\\p\nmid B_n}}\frac1p,\qquad
 \frac1X\sum_{n\leq X}f(n)\longrightarrow c_0,\\
 c_0=\sum_{r\geq1}2^{-r}\log(1+1/r)=A(1).
 \end{gathered}
\end{equation}
This full first moment is the only moment estimate imported from that
paper; the uniform prime counts were proved above. In particular, we
do not infer concentration of a sub-sum from the second
moment of the full sum, and do not import display (6) as written.

Here and below, the errors in \eqref{eq:uniform1} and
\eqref{eq:uniform2} have the following precise meaning. For fixed $0<a<b<1$
and $0<\epsilon<a/4$, let $L=\lceil2/a\rceil$ and let
$e_1(X;a,b,\epsilon)$ be the maximum of
$|A_1(p;X)/X-2^{-r}|$ over $X^a<p\leq X^b$, $1\leq r\leq L$, in
the indicated trimmed bands. Define $e_2$ as the analogous maximum for
\eqref{eq:uniform2}, with $r,t\leq L$; an empty maximum is zero.
The just-proved estimates give, for every integer $X\geq2^{1/\epsilon}$,
\[
 e_1\leq8LX^{-\epsilon}+X^{b-1},\qquad
 e_2\leq16LX^{-\epsilon}+X^{b-1}.
\]
Thus both errors tend to zero with the parameters fixed. For example,
both are less than any given $z>0$ whenever
\[
 X>\max\left\{2^{1/\epsilon},(32L/z)^{1/\epsilon},
                                  (2/z)^{1/(1-b)}\right\}.
\]
The only qualitative remainder retained here is the absolute error in
\eqref{eq:inputmean}: for every $z>0$ it is less than $z$ for all
sufficiently large $X$. No numerical onset for that imported first
moment, or effective exceptional-set rate for the final theorem, is
claimed.

We also use the reciprocal-prime bounds of Rosser--Schoenfeld
\cite{RS}. If $M(y)=\sum_{p\leq y}1/p$ and $B$ is the prime Mertens
constant, then
\begin{equation}\label{eq:RS}
 \begin{aligned}
 \log\log y+B-\frac1{2\log^2y}&<M(y)\quad(y>1),\\
 M(y)&<\log\log y+B+\frac1{2\log^2y}\quad(y\geq286).
 \end{aligned}
\end{equation}
Consequently, for $a\leq u\leq v\leq b$ and $X\geq286^{1/a}$, an
interval with exponent endpoints $u,v$ has reciprocal-prime mass
differing from $\log(v/u)$ by at most
\begin{equation}\label{eq:MX}
 E_X(a):=\frac1{a^2\log^2X}+2X^{-a}.
\end{equation}
The $2X^{-a}$ permits either inclusion convention at the two endpoints.

\begin{lemma}[Localized moment statement]\label{lem:local}
For every fixed $0<a<b<1$, in natural density,
\begin{equation}\label{eq:local}\tag{LOCAL}
 S_{a,b}(n):=\sum_{\substack{n^a<p\leq n^b\\p\nmid B_n}}\frac1p
 \longrightarrow \lambda_{a,b}:=\int_a^b g(t)\,dt.
\end{equation}
\end{lemma}
\begin{proof}
First choose $n$ uniformly from $\{1,\ldots,X\}$ and use fixed endpoints:
\[
 S^X(n)=\sum_{\substack{X^a<p\leq X^b\\p\nmid B_n}}\frac1p.
\]
Abbreviate $E=E_X(a)$, $H=\log(b/a)+E$, $\lambda=\lambda_{a,b}$,
and $L=\lceil2/a\rceil$. Fix $0<\epsilon<a/4$. Remove primes whose
exponent $\log p/\log X$ lies within distance $\epsilon$ of any $1/j$,
$1\leq j\leq L$. The union of removed intervals in $[a,b]$ has length
at most $2L\epsilon$ and at most $L$ components. Its reciprocal mass is
at most
\begin{equation}\label{eq:removed}
 R_X:=2L\epsilon/a+LE,
\end{equation}
by \eqref{eq:MX} and $\log(v/u)\leq(v-u)/a$. Its limiting $g$-mass
is at most $2L\epsilon/a$, since $w\leq1$. All other primes are in
the trimmed bands; every band met has index at most $L$.

\emph{First moment.} The count for $1\leq n\leq X$ differs from
$A_1(p;X)$ by at most $p\leq X^b$. On the retained primes apply
\eqref{eq:uniform1}; their total reciprocal mass is at most $H$.
There are at most $L$ band pieces, so \eqref{eq:MX} gives
\begin{equation}\label{eq:meanerror}
 |\E S^X-\lambda|\leq
 m_X:=(e_1+X^{b-1})H+2LE+4L\epsilon/a.
\end{equation}
Here one $LE$ bounds integration over the retained band pieces,
and the other $LE+4L\epsilon/a$ bounds the sum of the removed
prime mass and its limiting integral. In particular, first letting
$X\to\infty$ and then $\epsilon\downarrow0$ proves $\E S^X\to\lambda$.

\emph{Excluded pairs.} Choose an integer $s>1/a$. All the primes are
larger than $X^{1/s}$, so classes I and II of pp.~87--88 are empty.
Pairs containing a removed prime have reciprocal mass at most $2HR_X$.
For a retained prime in band $r$, $p^r<X^{1-r\epsilon}\leq
X^{1-\epsilon}$. Consecutive powers of the same prime have ratio
$p>X^a>X^\epsilon$. Thus failure of condition (5) for two distinct
retained primes requires
\begin{equation}\label{eq:ratio}
 X^{-\epsilon}\leq p^u/q^v\leq X^\epsilon
 \quad\hbox{for some }1\leq u,v\leq L.
\end{equation}
This is the ratio intended in the last excluded-pair display on p.~88,
which prints $p^u q^v$. The product with two positive exponents cannot
express proximity of two scales.
For fixed $p,u,v$, the possible exponents of $q$ lie in an interval
of length at most $2\epsilon/v\leq2\epsilon$. After intersecting with
$[a,b]$, \eqref{eq:MX} bounds their reciprocal mass by
$2\epsilon/a+E$. Summing over $p,u,v$ gives the explicit bound
\begin{equation}\label{eq:badpairs}
 \sum_{\text{excluded ordered pairs}}\frac1{pq}
 \leq2HR_X+L^2H(2\epsilon/a+E).
\end{equation}
This bound overcounts overlaps, which is harmless. It tends to a
quantity at most $(4L+2L^2)\log(b/a)\epsilon/a$.

\emph{Second moment and Chebyshev.} The diagonal contributes at most
\[
 \sum_{p>X^a}\frac1{p^2}\leq\frac1{X^a-1}.
\]
For the other, retained pairs satisfying (5), apply \eqref{eq:uniform2}.
Changing $n\geq\max(p,q)$ to $n\geq1$ costs at most $X^{b-1}$ in
each normalized count. Also, summing $w(\log p/\log X)/p$ band by
band differs from $\lambda$ by at most $LE$. Hence
\begin{align}
 \E (S^X)^2&\leq (\lambda+LE)^2+(e_2+X^{b-1})H^2+2HR_X\notag\\
 &\hspace{12mm}+L^2H(2\epsilon/a+E)+\frac1{X^a-1},\label{eq:second}\\
 \E(S^X-\lambda)^2&\leq
 (\lambda+LE)^2-\lambda^2+(e_2+X^{b-1})H^2+2HR_X\notag\\
 &\hspace{12mm}+L^2H(2\epsilon/a+E)+\frac1{X^a-1}
       +2\lambda m_X=:V_X.\label{eq:variance}
\end{align}
We used $-2\lambda\E S^X\leq-2\lambda^2+2\lambda m_X$.
Since $\lambda\leq\log(b/a)$,
\[
 \limsup_{X\to\infty}V_X
 \leq(12L+2L^2)\log(b/a)\epsilon/a.
\]
Chebyshev's inequality gives
$\Prob(|S^X-\lambda|>\eta)\leq V_X/\eta^2$. Letting
$\epsilon\downarrow0$ proves concentration of $S^X$.

\emph{Moving the endpoints.} Fix $0<\rho<1$. For
$X^{1-\rho}\leq n\leq X$, the symmetric difference of the two prime
ranges is contained in
\[
 (X^{a(1-\rho)},X^a]\ \cup\ (X^{b(1-\rho)},X^b].
\]
Consequently, for $X\geq286^{1/(a(1-\rho))}$,
\begin{equation}\label{eq:endpointshift}
 |S_{a,b}(n)-S^X(n)|\leq
 2\log\frac1{1-\rho}+2E_X(a(1-\rho)).
\end{equation}
Given $\eta>0$, choose $\rho$ so the logarithmic term is less than
$\eta/4$ and then $X$ so the error term is less than $\eta/4$.
There are at most $X^{1-\rho}$ discarded integers. Thus the density
of $|S_{a,b}-\lambda|>\eta$ is bounded by
$X^{-\rho}+\Prob(|S^X-\lambda|>\eta/2)$, which tends to zero.
\end{proof}

\begin{remark}[The endpoints in printed display (6)]\label{rem:printed6}
For arbitrary endpoints, whole digit bands cannot replace their actual
intersections with the interval. For example, take $a=3/10$, $b=9/20$.
The sum in printed (6), indexed by $1/b\leq k\leq1/a$, includes only
$k=3$ and gives $2^{-3}\log(4/3)$. In contrast, \eqref{eq:local} gives
\[
 2^{-3}\log\frac{1/3}{3/10}
 +2^{-2}\log\frac{9/20}{1/3}
 =2^{-3}\log(10/9)+2^{-2}\log(27/20).
\]
The two numbers differ: eight times their difference is
$\log((10/9)(27/20)^2/(4/3))=\log(243/160)>0$.
This is an endpoint-convention defect, independent of the missing
density qualifier in the printed theorem.
\end{remark}

\subsection{The finite sieve and its limit}\label{subsec:sieve}
For $n>1$ define the finite measure
\[
 \mu_n=\sum_{p\in\mathcal P_\delta(n)}\frac1p
             \delta_{\log p/\log n},\qquad
 \nu_\delta(dt)=\ind_{[\delta,\alpha]}(t)g(t)\,dt.
\]
The notation $\delta_u$ denotes a unit point mass at $u$.
For $n\geq286^{1/\alpha}$, \eqref{eq:RS} gives
\begin{equation}\label{eq:massbound}
 \mu_n([\delta,\alpha])
 \leq\log(\alpha/\delta)
 +\frac{\alpha^{-2}+\delta^{-2}}{2\log^2n}.
\end{equation}
In particular, with $C=1+\log(\alpha/\delta)$, this mass is at most
$C$ for
$n\geq\max(286^{1/\alpha},\exp\sqrt{(\alpha^{-2}+\delta^{-2})/2})$.
The same bound holds for the mass of $\nu_\delta$.

For completeness, we explain the measure passage without an independence
assumption. Partition $[\delta,\alpha]$ into $m$ half-open intervals of
length at most $h$. By \eqref{eq:local}, outside a set of density zero
each of their $m$ masses is within any prescribed $\zeta>0$ of its
$\nu_\delta$-mass. The sum of absolute errors of the bin masses is
at most $m\zeta$. The tensor telescoping identity bounds the error
of any union of bin boxes in the $k$th product measure by
$kC^{k-1}m\zeta$, using a larger $C$ if necessary for finitely many
initial $n$.

On $(0,1/2]$, $g(t)\leq3/4$: on band $r\geq2$ it is at most
$(r+1)2^{-r}\leq3/4$. A slab of width $\ell$ in the sum of the
$k$ coordinates has $\nu_\delta^{\otimes k}$-mass at most
\begin{equation}\label{eq:slab}
 \frac34\ell A(\alpha)^{k-1}.
\end{equation}
Indeed fix the first $k-1$ coordinates and integrate the last over an
interval of length at most $\ell$. Inner and outer bin approximations
to a simplex differ in a slab of width at most $2kh$. First choose
$h$, then $\zeta$, and then use the finite union of the exceptional
sets for the bins. This proves, in natural density, convergence of
$\mu_n^{\otimes k}$ on every simplex and every such slab to the
corresponding $\nu_\delta^{\otimes k}$ mass. In particular no joint
independence of the prime indicators is assumed in this step.

Set $K_\delta=\lfloor\alpha/\delta\rfloor$. Inclusion--exclusion gives
the exact finite identity
\begin{equation}\label{eq:sieve}
 F_\delta(n)=\sum_{k=0}^{K_\delta}(-1)^k
 \sum_{\substack{p_1<\cdots<p_k\in\mathcal P_\delta(n)\\d=p_1\cdots p_k\leq n^\alpha}}
 \frac{\lfloor n^\alpha/d\rfloor}{n^\alpha}.
\end{equation}
The $k=0$ term has $d=1$. For $0<\epsilon<\alpha$ and
$d\leq n^{\alpha-\epsilon}$, the error in replacing the floor term
by $1/d$ is at most $n^{-\epsilon}/d$. Summed over all such squarefree
products, it is at most
\begin{equation}\label{eq:floorerror}
 n^{-\epsilon}\prod_{p\in\mathcal P_\delta(n)}(1+1/p)
 \leq n^{-\epsilon}\exp(\mu_n([\delta,\alpha]))\leq n^{-\epsilon}e^C.
\end{equation}
For the remaining products the error is at most $1/d$. Their total
is bounded by
$\sum_{k=1}^{K_\delta}\mu_n^{\otimes k}
\{\alpha-\epsilon<\sum t_i\leq\alpha\}/k!$.
Its limit in density is at most
$(3/4)\epsilon\sum_{k=1}^{K_\delta}A(\alpha)^{k-1}/k!$,
by \eqref{eq:slab}. This tends to zero with $\epsilon$.

The repeated-prime part of the ordered $k$-tuple mass, for $k\geq2$,
is at most
\begin{equation}\label{eq:diagonal}
 \binom{k}{2}\left(\sum_{p>n^\delta}p^{-2}\right)C^{k-2}
 \leq\binom{k}{2}\frac{C^{k-2}}{n^\delta-1}.
\end{equation}
This follows by specifying one equal pair of coordinates and dropping
all other restrictions. Dividing the remaining distinct ordered sums
by $k!$ converts them into the unordered sums in \eqref{eq:sieve}.
Equations \eqref{eq:floorerror}--\eqref{eq:diagonal} and the simplex
convergence prove
\begin{equation}\label{eq:Fdelta}
 F_\delta(n)\longrightarrow
 c_{\mathrm D,\delta}:=\sum_{k=0}^{K_\delta}(-1)^kJ_k^{(\delta)}(\alpha)
 \quad\hbox{in natural density},
\end{equation}
where $J_k^{(\delta)}$ restricts all coordinates to $[\delta,\alpha]$.
Terms with $k>K_\delta$ are zero; a possible boundary simplex has
measure zero. It follows that $U_\delta\to1-c_{\mathrm D,\delta}$.

\subsection{Insufficient prime powers}\label{subsec:T}
The elementary valuation formula
\begin{equation}\label{eq:valuation}
 v_p(B_n)=\sum_{i\geq1}
 \left(\left\lfloor\frac{2n}{p^i}\right\rfloor
             -2\left\lfloor\frac n{p^i}\right\rfloor\right)
\end{equation}
follows by counting multiples of $p^i$ in factorials. Its summands
are zero or one; it is also the number of carries in adding $n+n$ in
base $p$. Each digit of $n$ at least $\lceil p/2\rceil$ forces a carry.
Fix $p$ and a nonnegative integer $V$. Write
$a_p=\lceil p/2\rceil$, $b_p=\lfloor p/2\rfloor$, and
$h_p=\lfloor\log_pX\rfloor$. A block of $p^{h_p}$ consecutive
integers has exactly
\[
 Q_{p,V}(h_p)=\sum_{j=0}^{\min(V,h_p)}
        \binom{h_p}{j}b_p^j a_p^{h_p-j}
\]
residues with at most $V$ large digits among the lowest $h_p$ digits.
There are at most $\lceil X/p^{h_p}\rceil\leq2X/p^{h_p}$ blocks,
including a partial one. Hence
\begin{equation}\label{eq:fixedp}
 \frac1X\#\{n\leq X:v_p(B_n)\leq V\}
 \leq\frac{2Q_{p,V}(h_p)}{p^{h_p}}.
\end{equation}
For $h_p\geq\max(V,1)$ the right side is at most
$2(V+1)h_p^V(a_p/p)^{h_p}$, since $b_p\leq a_p$ and
$\binom{h_p}{j}\leq h_p^V$. It tends to zero, including when $p=2$.

For an integer $H\geq2$, the part of $T(n)$ from $p>H$ is at most
$\sum_{m>H}m^{-2}\leq1/H$, by integration of $t^{-2}$.
If all $p\leq H$ have $v_p(B_n)>V$, the other part is at most
$H2^{-V-2}$. Given $\eta>0$, choose $H$ and $V$ so that
$1/H+H2^{-V-2}<\eta$. Then
\[
 \frac1X\#\{n\leq X:T(n)>\eta\}
 \leq\sum_{p\leq H}\frac{2Q_{p,V}(h_p)}{p^{h_p}}\longrightarrow0.
\]
Thus $T\to0$ in natural density with every discarded bound specified.

\subsection{Small primes and assembly}\label{subsec:G}
Fix $\rho>0$ with $\delta<1-\rho$. Pointwise,
$f(n)\geq G_\delta(n)+S_{\delta,1-\rho}(n)$.
The local sum is bounded for all sufficiently large $n$ by
\eqref{eq:RS}. A bounded sequence converging in density converges
in mean: if $|u(n)|\leq C$ and $u(n)\to u$ in density, the mean of
$|u(n)-u|$ is at most $\eta$ plus $C+|u|$ times the proportion where
$|u(n)-u|>\eta$; finitely many initial exceptions contribute their
fixed total divided by $X$. Apply this fact to \eqref{eq:local} and
then use \eqref{eq:inputmean} to obtain
\[
 \limsup_{X\to\infty}\frac1X\sum_{n\leq X}G_\delta(n)
 \leq c_0-\int_\delta^{1-\rho}g(t)\,dt.
\]
Let $\rho\downarrow0$. Monotone convergence gives
\begin{equation}\label{eq:G}
 \begin{aligned}
 \limsup_{X\to\infty}\frac1X\sum_{n\leq X}G_\delta(n)&\leq A(\delta),\\
 \limsup_{X\to\infty}\frac1X\#\{n\leq X:G_\delta(n)>\eta\}
 &\leq\frac{A(\delta)}\eta.
 \end{aligned}
\end{equation}
The second inequality is Markov's inequality. The full second moment
and a uniform-integrability argument are unnecessary here.

Deleting the measure on $(0,\delta)$ changes its total mass by
$A(\delta)$. Dropping the simplex restriction, or telescoping product
measures, gives for $k\geq1$
\[
 0\leq J_k-J_k^{(\delta)}
 \leq\frac{A(\alpha)^k-(A(\alpha)-A(\delta))^k}{k!}
 \leq\frac{A(\delta)A(\alpha)^{k-1}}{(k-1)!}.
\]
Consequently
\begin{equation}\label{eq:deltaerror}
 |c_{\mathrm D,\delta}-\cd(\alpha)|
 \leq A(\delta)e^{A(\alpha)}<2A(\delta).
\end{equation}
Given $\eta>0$, choose any sufficiently small $\delta>0$ so that
$2A(\delta)<\eta/4$. Outside the union of
\[
 \begin{gathered}
 \{|U_\delta-(1-c_{\mathrm D,\delta})|>\eta/4\},\\
 \{T>\eta/4\},\qquad \{G_\delta>\eta/4\},
 \end{gathered}
\]
the sandwich \eqref{eq:sandwich} implies $|N_\alpha-\cn|\leq\eta$.
The first two exceptional sets have density zero, and the third has
upper density at most $4A(\delta)/\eta$. Letting $\delta\downarrow0$
proves the first assertion of Theorem~\ref{thm:main}. The identity
$D_\alpha+N_\alpha=\lfloor n^\alpha\rfloor/n^\alpha$, whose error from
one is at most $n^{-\alpha}$, proves the second. Both sequences lie in
$[0,1]$, so the preceding bounded-convergence-in-density argument
proves Ces\`aro convergence as well.

\section{Failure of the printed statement as written}\label{sec:false}
\begin{proof}[Proof of Proposition~\ref{prop:nolimit}]
For $n=2^h$, all the terms in \eqref{eq:valuation} with $p=2$ vanish
except $i=h+1$, which equals one. Thus $4\nmid B_{2^h}$ and every
multiple of four is a nondivisor. The proposition's explicit lower
bound follows from $\lfloor y\rfloor\geq y-1$.

By Lemma~\ref{lem:alternating}, $\cn(\alpha)<1/4$ for every
$0<\alpha\leq1/2$. A limit along all integers would also be a limit
in natural density, and such a limit is unique: two different limits
have disjoint neighborhoods of radius less than half their separation,
whereas each neighborhood would have density one. The subsequence
$n=2^h$ therefore contradicts existence of an all-integer limit for
$N_\alpha$. Since $N_\alpha+D_\alpha\to1$, it also contradicts
existence of one for $D_\alpha$.

As $\alpha\downarrow0$, $s=\lfloor1/\alpha\rfloor\to\infty$ and
$A(\alpha)\leq2^{1-s}/s\to0$. The bracket
$0\leq\cn(\alpha)\leq A(\alpha)$ gives $\cn(\alpha)\to0$ and
$\cd(\alpha)=1-\cn(\alpha)\to1$.
\end{proof}
Thus the nondivisor direction of the printed small-$\alpha$ assertion
fails even with the density qualifier restored. Switching to divisors
repairs that limit, but does not produce an all-integer limit. These
conclusions concern $0<\alpha\leq1/2$; the interval $1/2<\alpha<1$
remains outside the statement proved here.

\section{Certified evaluation}\label{sec:numerics}
The accompanying program \texttt{code/density\_intervals.py} uses only
exact integers and rational numbers. We describe each error bound,
including the discretization error. The stated table uses $\alpha=1/s$
with $s\in\{2,3,4,5\}$, mesh $M=16384$, cutoff $R=64$, $K=8$ terms,
and $Q=2^{56}$. The program as shipped implements $\alpha=1/s$ for
integers $s\geq2$; the construction extends to any rational
$0<\alpha\leq1/2$ by clipping the top band at $\alpha$.

\subsection{Rational bin masses and logarithms}\label{subsec:bins}
Let $h=\alpha/M$ and $I_i=((i-1)h,ih]$, $1\leq i\leq M$.
First restrict $\nu(dt)=g(t)\,dt$ to $[1/R,\alpha]$, with $R>s$.
An intersection $[a,b]$ of a bin and a band $[1/(r+1),1/r]$ has mass
$2^{-r}\log(b/a)$. Zero-length intersections are ignored. Here
$1\leq b/a\leq(r+1)/r\leq3/2$. For $z=(b/a-1)/(b/a+1)$, put
\[
 \ell_J=2\sum_{j=0}^{J-1}\frac{z^{2j+1}}{2j+1}.
\]
Integrating the geometric series for $2/(1-z^2)$, or summing its
positive terms, gives the rational enclosure
\begin{equation}\label{eq:loginterval}
 \ell_J\leq\log(b/a)\leq
 \ell_J+\frac{2z^{2J+1}}{(2J+1)(1-z^2)}.
\end{equation}
Indeed each denominator in the omitted series is at least $2J+1$.
Here $0\leq z\leq1/5$, so the bound tends to zero. Increase $J$ until
the displayed error is less than $1/(Q2^{20})$. Round the lower and
upper bounds for $2^{-r}\log(b/a)$ down and up, respectively, to
integer multiples of $1/Q$. Sum segment bounds within each bin to get
\[
 q_i^-/Q\leq\nu(I_i\cap[1/R,\alpha])\leq q_i^+/Q.
\]
Every operation defining these integers is exact.

\subsection{Simplex bounds and exact convolution}\label{subsec:convolution}
For a tuple of bins, $\sum i_j\leq M$ implies that its whole box lies
in $\sum t_j\leq\alpha$. Conversely, an intersecting box has
$\sum i_j\leq M+k-1$, and in particular $\sum i_j\leq M+k$.
Thus, with $P_\pm(z)=\sum_{i=1}^M q_i^\pm z^i$, lower and upper
bounds for the truncated-measure $J_k$ are
\begin{equation}\label{eq:poly}
 L_k=\frac{\sum_{j=0}^M[z^j]P_-(z)^k}{k!Q^k},\qquad
 U_k=\frac{\sum_{j=0}^{M+k}[z^j]P_+(z)^k}{k!Q^k}.
\end{equation}
Repeated bins are allowed: products of their masses are exactly the
product-measure masses of boxes. Divide by $k!$ because $J_k$ uses
ordered integration variables with that prefactor. Alternating signs
are applied outwards, using $[-U_k,-L_k]$ when $k$ is odd.

To compute the coefficients, pack each polynomial into base $2^B$,
where
\[
 B=8\left\lfloor\frac{56K+8}{8}\right\rfloor=456
 \quad\hbox{for }K=8.
\]
More generally the code replaces $56$ by the chosen bit precision.
It checks $\sum_iq_i^+<Q$. Thus every coefficient and every prefix
sum in any power $1\leq k\leq K$ is strictly less than
$Q^k\leq Q^K<2^B-1$. There are no carries between packed coefficients.
Integer multiplication and masking after degree $M+K$ therefore give
exact truncated polynomial multiplication; discarded high degrees
cannot affect lower ones because all degrees are nonnegative.
Reduction of a selected prefix modulo $2^B-1$ sums its base-$2^B$
digits exactly: $2^B\equiv1$ modulo $2^B-1$, and the digit sum is
less than that modulus. This proves the convolution algorithm, without
any floating-point rounding assumption.

\subsection{All remaining errors and convergence}\label{subsec:errors}
The omitted mass below $1/R$ satisfies
\[
 \tau:=\int_0^{1/R}g(t)\,dt
 =\sum_{r\geq R}2^{-r}\log(1+1/r)
 \leq\tau_R:=\frac{2^{1-R}}R.
\]
Writing $A=A(\alpha)$, the absolute change in the full alternating
series is bounded by
\[
 \sum_{k\geq1}\frac{A^k-(A-\tau)^k}{k!}
 =e^A-e^{A-\tau}\leq\tau e^A\leq\tau_R e^{A_s}<2\tau_R.
\]
Truncating the $k$-series adds at most $E_K(A_s)$ by
\eqref{eq:tail}. Add
\begin{equation}\label{eq:radius}
 2\tau_R+E_K(A_s)
\end{equation}
to each end of the interval computed from \eqref{eq:poly}. For $s=2$,
$R=64$, $K=8$, this radius is strictly less than
$0.00000000001078182876$. Nondivisor bounds are one minus the reversed
divisor bounds. Decimal endpoints are rounded outwards using integer
division, so they remain rational enclosures.

Here is also an explicit convergence guarantee. Let
$\eta_Q=\sum_i(q_i^+-q_i^-)/Q$. There are at most $M+R$ positive
bin--band intersections, since the common refinement of the two
partitions has at most the sum of their numbers of pieces. For each
segment the scaled logarithm enclosure has length less than $2^{-20}$;
after rounding its integer width is at most two. Therefore
\begin{equation}\label{eq:etaQ}
 \eta_Q\leq\frac{2(M+R)}Q.
\end{equation}
For the exact bin masses, the boxes included in the upper simplex
bound but not the lower one lie in a slab of width at most $2kh$.
By \eqref{eq:slab}, summing their contributions divided by $k!$ is
at most
\[
 \sum_{k\geq1}\frac{2kh(3/4)A_s^{k-1}}{k!}
 =\frac32h e^{A_s}<3h.
\]
Let $m_-$ and $m_+$ be the total lower and upper rounded masses.
They differ by $\eta_Q$, and $m_+\leq A_s+\eta_Q$. The sum of
the rounding contributions to the interval width is at most
$\sum_{k\geq1}(m_+^k-m_-^k)/k!
\leq\eta_Q e^{m_+}<3\eta_Q$ when $A_s+\eta_Q\leq1$.
Together with \eqref{eq:radius}, the unrounded enclosure width is at most
\begin{equation}\label{eq:width}
 3h+3\eta_Q+4\tau_R+2E_K(A_s).
\end{equation}
Reporting to $d$ decimal places adds at most $2\cdot10^{-d}$ to
this width. Increasing $M,R,K,Q$ so that each displayed bound is below
its allotted tolerance gives any prescribed rational accuracy.
All choices have explicit bounds: one may test \eqref{eq:width}
using rational upper bounds $3h$, $6(M+R)/Q$, $4\tau_R$, and
$2E_K(A_s)$, then increase the parameters until their sum is small
enough. The assertion $\sum q_i^+<Q$ eventually holds since
$A_s\leq1/4$ and \eqref{eq:etaQ} tends to zero.

For a real input $\alpha$ specified by computable rational
approximations, a modulus also follows from the same slab estimate.
If $0<\alpha\leq\beta\leq1/2$, then
\[
 0\leq J_k(\beta)-J_k(\alpha)
 \leq\frac{3(\beta-\alpha)A(\beta)^{k-1}}{4k!},\quad k\geq1.
\]
Summing shows $|\cd(\beta)-\cd(\alpha)|$ and
$|\cn(\beta)-\cn(\alpha)|$ are at most
$(3/4)e^{1/4}(\beta-\alpha)<2(\beta-\alpha)$.
This supplies the approximation interpretation in
Remark~\ref{rem:scope}; it asserts no algorithm for an unspecified,
noncomputable real input.

\subsection{Certified intervals and reproduction}\label{subsec:table}
The following endpoints are transcribed exactly from
\texttt{code/density\_16384.txt}. They are outward-rounded rational
intervals, not confidence intervals.
\begin{center}
\small
\begin{tabular}{c|c|c}
$\alpha$ & $\cn(\alpha)$: nondivisors & $\cd(\alpha)$: divisors\\ \hline
$1/2$ & $[0.160401041573,\ 0.160401748642]$
      & $[0.839598251358,\ 0.839598958427]$\\
$1/3$ & $[0.059862571395,\ 0.059862608229]$
      & $[0.940137391771,\ 0.940137428605]$\\
$1/4$ & $[0.023932466470,\ 0.023932468474]$
      & $[0.976067531526,\ 0.976067533530]$\\
$1/5$ & $[0.009987263535,\ 0.009987263648]$
      & $[0.990012736352,\ 0.990012736465]$
\end{tabular}
\end{center}
For $\alpha=1/2$ the displayed width is $0.000000707069$.
The mesh-$4096$ interval is $[0.160399449123,0.160402279170]$,
of width $0.000002830047$. The ratio of these displayed widths is
greater than four, since $2830047>4\cdot707069$.

Run the following from the directory containing \texttt{main.tex};
the Python programs need only the standard library:
\begingroup\footnotesize
\begin{verbatim}
python3 code/density_intervals.py --mesh 4096 \
  > density_4096.rerun.txt
python3 code/density_intervals.py --mesh 16384 \
  > density_16384.rerun.txt
cmp code/density_4096.txt density_4096.rerun.txt
cmp code/density_16384.txt density_16384.rerun.txt
python3 code/check_convolution.py
\end{verbatim}
\endgroup
The convolution check compares packed powers to independent
nested-loop convolution for $s\in\{2,3,4,5\}$, $M\in\{8,16,31\}$,
$R=12$, both choices of coefficient bounds, and $k=1,2,3,4$:
$4\cdot3\cdot2\cdot4=96$ exact comparisons. It prints
\begingroup\footnotesize
\begin{verbatim}
96 packed polynomial powers agree with independent nested-loop convolutions.
\end{verbatim}
\endgroup
These finite checks test the implementation; the certification rests on
\eqref{eq:loginterval}--\eqref{eq:width}. Neither the table nor the
checks supply an effective exceptional-set rate for the density theorem.

\section*{Provenance and verification status}
The density theorem, its proof and the certified evaluation were found by
a GPT-6 Astra agent (OpenAI) on 2026-09-09; the present text, which applies
the referees' repairs, was revised by GPT-6 Astra on 2026-10-10. Two
isolated cross-vendor referee reports by Claude Opus (Anthropic), dated
2026-10-10, both returned \emph{PASS-WITH-REPAIRS} with no fatal gap.
Both asked for the same three required repairs, all applied here:
the decreasing-term inequality behind $0\leq\cn(\alpha)\leq A(\alpha)$
(Lemma~\ref{lem:alternating}); the localized moment statement written out
(Lemma~\ref{lem:local}); and the scope $0<\alpha\leq1/2$
(Remark~\ref{rem:scope}). Referee~1 also suggested an optional
simplification of the small-prime argument, applied in
Section~\ref{subsec:G}. In a second round both referees asked for a
self-contained proof of the uniform forms of the two counts taken from
the 1975 paper; this is Lemma~\ref{lem:nested}, which both confirmed in a
third round (\emph{PASS}). Each referee reran the certificate program at
both meshes with byte-identical output, and computed the four constants
by an independent method; the converged values fell inside the certified
intervals. A separate Claude verifier reproduced the certificates byte
for byte and checked $v_2\bigl(\binom{2^{h+1}}{2^h}\bigr)=1$ exactly for
$1\leq h\leq300$. Empirical proportions for $n$ between about $10^{9}$
and $10^{24}$ lie roughly 9--17\% below the limits, depending on $\alpha$
and on the range of $n$, which is consistent with slow convergence in
$n$. These are sanity checks, not evidence for the theorem.
No formal (Lean) verification has been carried out.

\begin{thebibliography}{99}
\bibitem{EGRS}
P. Erd\H{o}s, R. L. Graham, I. Z. Ruzsa and E. G. Straus,
\emph{On the prime factors of $\binom{2n}{n}$},
Mathematics of Computation \textbf{29} (1975), 83--92.
\href{https://doi.org/10.1090/S0025-5718-1975-0369288-3}
{doi:10.1090/S0025-5718-1975-0369288-3}.
\bibitem{RS}
J. B. Rosser and L. Schoenfeld,
\emph{Approximate formulas for some functions of prime numbers},
Illinois Journal of Mathematics \textbf{6} (1962), 64--94.
\bibitem{EP377}
T. F. Bloom, \emph{Erd\H{o}s Problems}, Problem 377,
\url{https://www.erdosproblems.com/377}, accessed 10 October 2026.
\bibitem{FK21}
K. Ford and S. Konyagin,
\emph{Divisibility of the central binomial coefficient $\binom{2n}{n}$},
Transactions of the American Mathematical Society \textbf{374} (2021), 923--953.
\bibitem{FC5792}
Google DeepMind, \emph{formal-conjectures}, pull request 5792 (merged
12 September 2026),
\url{https://github.com/google-deepmind/formal-conjectures/pull/5792}.
\bibitem{Po15}
C. Pomerance, \emph{Divisors of the middle binomial coefficient},
American Mathematical Monthly \textbf{122} (2015), 636--644.
\bibitem{Po26}
C. Pomerance, \emph{Remarks on the middle binomial coefficient},
Integers \textbf{26} (2026), A47.
\bibitem{Sa18}
C. Sanna, \emph{Central binomial coefficients divisible by or coprime to
their indices}, International Journal of Number Theory \textbf{14} (2018),
1135--1141.
\end{thebibliography}
\end{document}
